\documentclass[11pt]{article}
\usepackage[a4paper,margin=25mm]{geometry}
\usepackage{amsmath,amssymb,amsthm,booktabs,microtype}
\usepackage[T1]{fontenc}
\usepackage{lmodern}
\usepackage[dvipsnames]{xcolor}
\usepackage[colorlinks=true,linkcolor=MidnightBlue,citecolor=MidnightBlue,urlcolor=MidnightBlue]{hyperref}
\usepackage{fancyhdr}
\pagestyle{fancy}\fancyhf{}
\fancyhead[L]{\small Three-uniform Berge-$C_4$ saturation}
\fancyhead[R]{\small Dev D. Goyal}
\fancyfoot[C]{\thepage}
\setlength{\headheight}{14pt}
\hypersetup{pdftitle={Three-uniform Berge-C4 saturation: linear bounds with coefficients 5/8 and 3/4, the cone spectrum and exact small values},
pdfauthor={Dev D. Goyal},
pdfsubject={Fully autonomous agentic research},
pdfkeywords={Berge hypergraphs, Berge cycles, saturation numbers, three-uniform hypergraphs, saturation spectrum, Lean 4}}
\newtheorem{theorem}{Theorem}
\newtheorem{lemma}[theorem]{Lemma}
\newtheorem{proposition}[theorem]{Proposition}
\newtheorem{corollary}[theorem]{Corollary}
\theoremstyle{remark}
\newtheorem{remark}[theorem]{Remark}
\newcommand{\sat}{\operatorname{sat}}
\newcommand{\codeg}{\operatorname{codeg}}
\title{\textbf{Three-uniform Berge-$C_4$ saturation:\\ linear bounds with coefficients $5/8$ and $3/4$,\\ the cone spectrum and exact small values}}
\author{Dev D. Goyal}
\date{31 October 2026}
\begin{document}
\maketitle
\begin{abstract}
Let $s(n)=\sat_3(n,\text{Berge-}C_4)$ be the least number of hyperedges of
a three-uniform hypergraph on $n$ vertices that contains no Berge
four-cycle but contains one after any absent triple is added. Neither the
hypergraph nor the added triple is required to be linear. English,
Graber, Kirkpatrick, Methuku and Sullivan showed $s(n)\le n-2$ for
$n\ge7$, and English, Gerbner, Methuku and Tait showed that $s(n)$ grows
linearly. We prove that for every $n\ge9$
\[
\max\Bigl(\Bigl\lceil\tfrac{n-1}{2}\Bigr\rceil,\Bigl\lceil\tfrac{5n-3623}{8}\Bigr\rceil,
\Bigl\lceil\tfrac{5n+89}{8}-\tfrac{\sqrt{270n+1891}}{4}\Bigr\rceil\Bigr)
\le s(n)\le\Bigl\lceil\tfrac{3n-9}{4}\Bigr\rceil ,
\]
so that $5/8\le\liminf s(n)/n\le\limsup s(n)/n\le3/4$. We also classify
the saturated hypergraphs in which one vertex lies in every hyperedge
(cones) for $n\ge5$ and determine their sizes: exactly the integers from
$\lceil(3n-5)/4\rceil$ to $n-2$, together with $n-1$ when $n\equiv1\pmod
3$. Hence every integer in $[\lceil(3n-9)/4\rceil,n-2]$ is the size of a
saturated hypergraph. A vertex $r$ such that every other vertex has both
a one-edge and a two-edge Berge path to $r$ forces at least $3(n-1)/4$
hyperedges, with equality exactly for bouquets of three-edge blocks.
Exhaustive certificates give $s(n)$ for $n\le13$. Lean~4 proves the upper
bound for every $n\ge9$, for actual hypergraphs, Berge cycles and
saturation under every absent triple, and the values $s(n)$ for $n\le5$;
the lower bound, the cone results and the values for $6\le n\le13$ have
conventional proofs or checked certificates.
\end{abstract}

\section{Definitions and results}

All hypergraphs are simple and three-uniform, and the vertex set may
contain isolated vertices. A \emph{Berge-$C_4$} in $H$ consists of four
distinct vertices $v_0,v_1,v_2,v_3$ and four distinct hyperedges
$e_0,e_1,e_2,e_3$ of $H$ with $\{v_i,v_{i+1}\}\subseteq e_i$, indices
modulo~$4$; a hyperedge may contain further cycle vertices. A \emph{Berge
path of length $\ell$} has distinct vertices $v_0,\dots,v_\ell$ and
distinct hyperedges $e_1,\dots,e_\ell$ with $\{v_{i-1},v_i\}\subseteq e_i$.
A hypergraph $H$ is \emph{saturated} if it contains no Berge-$C_4$ but
$H+f$ contains one for every three-element set $f\notin E(H)$. This is
the unrestricted notion of~\cite{EGKMS}: $f$ may share two vertices with
a hyperedge of $H$. We write $s(n)$ for the least number of hyperedges of
a saturated hypergraph on $n$ vertices.

English, Graber, Kirkpatrick, Methuku and Sullivan~\cite[Theorem~5.7]{EGKMS}
gave the upper bound $s(n)\le n-2$ for $n\ge7$, and English, Gerbner,
Methuku and Tait~\cite{EGMT} proved $\sat_k(n,\text{Berge-}F)=O(n)$ for
every graph $F$ and $3\le k\le5$. The elementary lower bound
$s(n)\ge\lceil(n-1)/2\rceil$ follows from connectivity
(Lemma~\ref{lem:basic}). Wang and Zhang~\cite{WZ} studied the different
\emph{linear} saturation problem, in which $H$ is linear and only additions
preserving linearity are tested.

\begin{theorem}\label{thm:main}
For every $n\ge4$,
\begin{equation}\label{eq:lower}
5n\le 8s(n)+3623\qquad\text{and}\qquad 5n\le 8s(n)+2\sqrt{432\,s(n)+1}+19,
\end{equation}
and $s(n)\ge\lceil(n-1)/2\rceil$. For every $n\ge9$,
$s(n)\le\lceil(3n-9)/4\rceil$. Consequently, for $n\ge9$, $s(n)$ lies
between the maximum displayed in the abstract and $\lceil(3n-9)/4\rceil$,
and $5/8\le\liminf s(n)/n\le\limsup s(n)/n\le 3/4$.
\end{theorem}

The square-root form of~\eqref{eq:lower} inverts to the third term of the
abstract (Section~\ref{sec:lower}); it is the better of the two for
$n\le12751$. We do not know whether $s(n)/n$ converges.
The constants in~\eqref{eq:lower} are not optimized.

A \emph{cone} is a hypergraph with a vertex $r$ in every hyperedge. Its
\emph{link} is the graph $G$ on $V\setminus\{r\}$ with $xy\in E(G)$ when
$rxy\in E(H)$. Write $D(p,q)$ for the double star with adjacent centres
carrying $p\ge1$ and $q\ge1$ leaves, $K_{1,t}$ for the star with $t$
edges and $K_3$ for the triangle.

\begin{theorem}\label{thm:cone}
Let $n\ge5$ and let $H$ be a cone on $n$ vertices with link $G$. Then $H$
is saturated if and only if either $G$ has no isolated vertex and its
components are double stars, triangles and at most one star $K_{1,t}$ with
$t=1$ or $t\ge3$, or $G$ has exactly one isolated vertex and all other
components are triangles. The sizes of saturated cones on $n$ vertices
are exactly the integers in $[\lceil(3n-5)/4\rceil,\,n-2]$, together with
$n-1$ when $n\equiv1\pmod3$.
\end{theorem}

\begin{corollary}\label{cor:interval}
For $n\ge9$ every integer in $[\lceil(3n-9)/4\rceil,\,n-2]$, a set of
$\lfloor(n+5)/4\rfloor$ integers, is the size of a saturated hypergraph on
$n$ vertices. The construction of Theorem~\ref{thm:main} has one hyperedge
fewer than the smallest saturated cone.
\end{corollary}

\begin{theorem}\label{thm:root}
Let $H$ be a three-uniform hypergraph on $n$ vertices with a vertex $r$
such that every other vertex has a one-edge and a two-edge Berge path to
$r$. Then $|E(H)|\ge3(n-1)/4$, with equality exactly when $H$ is a union of
blocks $\{r,b,c\},\{r,b,x\},\{r,c,y\}$ that are pairwise disjoint outside
$r$.
\end{theorem}

Theorem~\ref{thm:root} needs neither freeness nor saturation. It shows
that the gluing interface of Section~\ref{sec:upper} cannot beat the
coefficient $3/4$; it is not a lower bound for $s(n)$.

\begin{theorem}\label{thm:small}
For $n=0,1,\dots,13$,
\[
s(n)=0,0,0,1,3,3,4,4,5,5,6,6,7,8 .
\]
At $n=11$ there are exactly two saturated hypergraphs with six hyperedges
up to isomorphism, $\{023,034,123,156,278,\{3,9,10\}\}$ and
$\{012,023,045,126,178,\{2,9,10\}\}$.
\end{theorem}

The values agree with $\lceil(3n-9)/4\rceil$ for $9\le n\le13$; whether
$s(n)=\lceil(3n-9)/4\rceil$ for all $n\ge9$ is open.

\section{Basic facts}\label{sec:basic}

\begin{lemma}[path criterion]\label{lem:path}
Let $H$ contain no Berge-$C_4$ and let $f\notin E(H)$ be a triple. Then
$H+f$ contains a Berge-$C_4$ if and only if two vertices of $f$ are the
ends of a Berge path of length three in $H$.
\end{lemma}

\begin{proof}
A Berge-$C_4$ in $H+f$ uses $f$; deleting $f$ from it leaves a path of
length three whose ends lie in $f$. Conversely, closing such a path by $f$
gives a Berge-$C_4$, since $f$ differs from the hyperedges of $H$; the
third vertex of $f$ may lie on the path.
\end{proof}

A \emph{leaf} is a vertex of degree one. For a saturated $H$ on $n\ge4$
vertices no hyperedge has three leaves: two of them with an outside vertex
form an absent triple, and no Berge path of length three can end at a leaf
of such a hyperedge and at a vertex outside it, nor join two of its
leaves. A hyperedge with exactly two leaves is \emph{pendant}; $t$ denotes
the number of pendant hyperedges.

\begin{lemma}\label{lem:basic}
Let $H$ be saturated on $n\ge4$ vertices with $m$ hyperedges. Then $H$ has
at most one isolated vertex and at most two components. If $H$ is
disconnected, it has no leaf and $3m\ge2(n-1)$. In all cases
$m\ge\lceil(n-1)/2\rceil$.
\end{lemma}

\begin{proof}
If $x,y,z$ lie in three different components (isolated vertices
included), the triple $xyz$ is absent and no path joins two of its
vertices, contradicting Lemma~\ref{lem:path}. Suppose $H$ has two
components, and let $u,v$ lie in one of them and $w$ in the other. The
triple $uvw$ is absent and only $u,v$ can be joined, so every pair of
vertices in a common component is joined by a path of length three.
Then no component contains a leaf $x$ with hyperedge $xab$: a path of
length three from $a$ to $b$ cannot use $xab$ (at an end it would continue
from the leaf $x$, in the middle it would need two internal vertices in
$\{x\}$), and a path avoiding $xab$ closes with $xab$ to a Berge-$C_4$ in
$H$. So at most one vertex (an isolated one) has degree below two, and
$3m\ge2(n-1)$. If $H$ is connected, its incidence graph is connected,
so $n+m-1\le 3m$, that is, $m\ge(n-1)/2$.
\end{proof}

\begin{lemma}\label{lem:shadow}
If $H$ contains no Berge-$C_4$, its shadow (the graph of pairs covered by
hyperedges) contains no $K_5$.
\end{lemma}

\begin{proof}
First, four distinct triples $J$ inside a five-set $S$ always contain a
Berge-$C_4$. The degrees sum to $12$. If a vertex $r$ has degree four, its
link has four edges on four vertices, hence contains a path $abcd$; the
fourth link edge is $ac$, $ad$ or $bd$, and the triples for $ac$ (or $ad$),
$ab$, $bc$, $cd$ give the cycle $r,a,b,c$, while for $bd$ the reversed path
gives $r,d,c,b$. If some vertex has degree zero, all four triples on the
other four vertices are present and contain a Berge-$C_4$. If $z$ has
degree one, the other three triples are, up to renaming, $abc,abd,acd$, and
$z$'s triple contains $ab$ or $bc$; the cores $a,b,c,d$ and $b,c,a,d$ use
$z$'s triple, $abc$, $acd$, $abd$ in that order. The remaining degree
sequence is $(3,3,2,2,2)$; the complements of the four triples in $S$ form
a graph with degrees $(1,1,2,2,2)$, which is $P_5$ or a triangle plus a
disjoint edge. For $P_5=12345$ the triples are $345,145,125,123$ and the
core $1,2,3,4$ uses $125,123,345,145$; for the triangle $123$ and edge
$45$ the triples are $345,145,245,123$ and the core $1,2,4,5$ uses
$123,245,345,145$.

Now let the shadow contain $K_5$ on $S=\{1,\dots,5\}$, let $T$ be the set
of hyperedges inside $S$, and assume $|T|\le3$. Every pair of $S$ not
covered by $T$ is covered by a hyperedge meeting $S$ in exactly that pair,
and distinct pairs give distinct hyperedges. Up to relabelling, the
following cases exhaust $T$; each row gives a core and four distinct
covering hyperedges in cyclic order, where a two-digit entry is the
hyperedge chosen for that pair:
\begin{center}\small
\begin{tabular}{lll}
\toprule
$T$ & core & hyperedges\\
\midrule
$\emptyset$ & $1,2,3,4$ & $12;23;34;14$\\
$123$ & $1,2,4,5$ & $123;24;45;15$\\
$123,124$ & $1,2,3,5$ & $124;123;35;15$\\
$123,145$ & $2,3,4,5$ & $123;34;145;25$\\
$123,124,125$ & $1,2,3,4$ & $125;123;34;124$\\
$123,124,134$ & $1,2,3,5$ & $124;123;35;15$\\
$123,124,135$ & $1,2,4,3$ & $123;124;34;135$\\
$123,124,345$ & $1,2,3,5$ & $124;123;345;15$\\
\bottomrule
\end{tabular}
\end{center}
(For $|T|=3$ the complements of the triples form a triangle, a star, a
path $P_4$ or a path $P_3$ plus an edge.) Each case contains a Berge-$C_4$.
\end{proof}

\begin{lemma}\label{lem:pendant}
A saturated hypergraph on $n\ge4$ vertices has at most four pendant
hyperedges.
\end{lemma}

\begin{proof}
Let $axy$ be pendant with leaves $x,y$. No path of length three joins $x$
and $y$. For $z\notin\{a,x,y\}$ the absent triple $xyz$ therefore needs a
path of length three from $x$ or $y$ to $z$, which starts with $axy$ and
passes to $a$. If $z$ is a leaf of a second pendant hyperedge $bzw$, this
path ends with $bzw$, so its middle hyperedge contains $a$ and $b$, and
$a\ne b$. The attachment vertices of pendant hyperedges are therefore
distinct and pairwise covered in the shadow, so Lemma~\ref{lem:shadow}
allows at most four.
\end{proof}

Four is attained: the triples $012,013,023$ with one pendant triple at each
of $0,1,2,3$ form a saturated hypergraph on $12$ vertices.

\section{The upper construction}\label{sec:upper}

Let $r,b,c$ be vertices and $X,Y$ disjoint nonempty sets of further
vertices. For $Q\subseteq\{r,b,c\}$ let $H(X,Y,Q)$ have the hyperedges
$E_0=rbc$, $E_x=rbx$ ($x\in X$), $F_y=rcy$ ($y\in Y$) and, for each $u\in Q$,
one pendant triple $u\,p_u\,q_u$ with two new vertices. The \emph{block}
is $B=H(\{x\},\{y\},\emptyset)=\{rbc,rbx,rcy\}$.

\begin{lemma}\label{lem:base}
$H(X,Y,Q)$ is saturated. Every non-root vertex of $H(X,Y,Q)$ has a
one-edge or a two-edge path to $r$, and every non-root vertex of $B$ has
both.
\end{lemma}

\begin{proof}
Only $r,b,c$ have degree at least two, and the four vertices of a
Berge-$C_4$ have degree at least two, so there is none. Call the other
vertices leaves. Two leaves in different hyperedges are joined through
$E_0$ by choosing distinct core vertices in their own hyperedges (two
pendant hyperedges have different attachments). A vertex $x\in X$ reaches
the three core vertices by $x,E_x,b,E_0,c,F_y,r$; $x,E_x,r,F_y,c,E_0,b$;
$x,E_x,b,E_0,r,F_y,c$, and symmetrically for $Y$. A pendant leaf at $u$
reaches each other core vertex $z$ through $u,w,z$ with $w$ the third core
vertex, using $E_0$ for one step and $E_x$ or $F_y$ for the other. Every
absent triple contains one of these pairs, or consists of a pendant leaf,
its sibling and its attachment (a present hyperedge), or is $rbc$ (present).
Lemma~\ref{lem:path} applies. The path statements are immediate: $b,c,X,Y$
and leaves at $r$ are adjacent to $r$; a leaf at $b$ or $c$ reaches $r$
through its attachment and $E_0$; in $B$ the two-edge paths are
$b,c,r$; $c,b,r$; $x,b,r$; $y,c,r$.
\end{proof}

\begin{lemma}[rooted gluing]\label{lem:glue}
Glue one $H(X,Y,Q)$ and any number of copies of $B$ at $r$, disjointly
otherwise. The result is saturated.
\end{lemma}

\begin{proof}
Every hyperedge has its non-root vertices in one piece. Around a
Berge-$C_4$ at most one vertex is $r$, and each other cycle vertex lies
in two consecutive cycle hyperedges, so all four hyperedges lie in one
piece; pieces are free (a block has only three hyperedges). An absent
triple inside one piece and $r$ is handled by that piece. Otherwise it
contains non-root vertices $u,w$ of different pieces; one of them lies in
a block, so one has a one-edge and the other a two-edge path to $r$, and
the concatenation is a path of length three because the pieces share only
$r$.
\end{proof}

\begin{proof}[Proof of the upper bound in Theorem~\ref{thm:main}]
For $n\ge9$ let $n_0\in\{9,10,11,12\}$ with $n_0\equiv n\pmod4$ and take
$|X|,|Y|,|Q|=(1,1,2),(2,1,2),(1,1,3),(2,1,3)$, a base with $n_0$ vertices
and $5,6,6,7$ hyperedges. Gluing $(n-n_0)/4$ blocks adds four vertices and
three hyperedges each, which gives $\lceil(3n-9)/4\rceil$ hyperedges.
\end{proof}

The construction is formalized for every $n\ge9$
(Section~\ref{sec:lean}). The four bases used there are
$\{012,013,024,056,178\}$,
$\{012,013,014,025,067,189\}$,
$\{012,013,024,056,178,\{2,9,10\}\}$ and
$\{012,013,014,025,067,\{1,8,9\},\{2,10,11\}\}$ with root $0$.

\section{The lower bound}\label{sec:lower}

Let $H$ be saturated on $n\ge4$ vertices with $m$ hyperedges. A
\emph{one-leaf hyperedge} $abx$ has a leaf $x$ and $\deg a,\deg b\ge2$.
Call $x$ \emph{bad} if $\codeg(a,b)=1$ and \emph{good} otherwise; $q$
counts bad leaves. A good leaf is \emph{supported} if some hyperedge
containing $a,b$ has no leaf, and \emph{unsupported} otherwise; $A$ and
$B$ count supported and unsupported good leaves. There are $A+B+q+2t$
leaves.

\begin{lemma}\label{lem:local}
Each vertex $a$ lies in at most six hyperedges $ab_ix_i$ with $x_i$ a bad
leaf.
\end{lemma}

\begin{proof}
The $b_i$ are distinct. Join $b_i$ and $b_j$ in an auxiliary graph $T$
when some hyperedge contains both. If $b_ib_j,b_jb_k\in T$ through
different hyperedges $f,g$, then $a,b_i,b_j,b_k$ with $ab_ix_i,f,g,ab_kx_k$
is a Berge-$C_4$ (the hyperedge $ab_ix_i$ contains no other $b_j$). So $f=g=
b_ib_jb_k$, and no further vertex is adjacent to this triangle: the
components of $T$ have at most three vertices. A path of length three
from $x_i$ to $x_j$ starts with $ab_ix_i$, ends with $ab_jx_j$, and its
middle hyperedge must contain $b_ib_j$ (a middle hyperedge through $a$
and $b_i$ or $b_j$ would repeat an end hyperedge because the codegrees are
one); conversely $b_ib_j\in T$ gives such a path. Bad leaves from three
components of $T$ would form an absent triple with no joined pair. So $T$
has at most two components and at most six vertices.
\end{proof}

\begin{lemma}\label{lem:sqrt}
$q(q-2)\le432m$ and $q\le1803$.
\end{lemma}

\begin{proof}
Let $P$ be the graph on bad leaves joined by paths of length three. Any
three bad leaves form an absent triple, so the complement of $P$ is
triangle-free and Mantel's theorem gives $|E(P)|\ge\binom q2-\lfloor
q^2/4\rfloor\ge q(q-2)/4$. A joined pair has a middle hyperedge $g$ and
internal vertices $a,b\in g$, and each end leaf lies in a bad-leaf
hyperedge through $a$ or $b$. By Lemma~\ref{lem:local} a pair $a,b$
accounts for at most $36$ pairs of leaves, a hyperedge for at most
$108$, so $|E(P)|\le108m$ and $q(q-2)\le432m$.

For the absolute bound, let $w_a\in[1,6]$ be the number of bad-leaf
hyperedges through a core endpoint $a$; then $D=\sum w_a=2q$. Process the
shadow pairs among core endpoints in order of decreasing weight $w_aw_b$,
assigning each to a still unused hyperedge containing it when one exists.
The graph $F$ of assigned pairs has no four-cycle (its assigned hyperedges
would form a Berge-$C_4$). An unassigned pair $p$ lies only in hyperedges
already assigned to pairs of weight at least that of $p$, and a
hyperedge contains three pairs; hence the total shadow weight is at most
$3W$ with $W=\sum_{ab\in F}w_aw_b$. With $z_v=\sum_{u\sim v}w_u$, the
weighted Cauchy--Schwarz inequality gives $\sum_vw_vz_v^2\ge4W^2/D$. Two
vertices of $F$ have at most one common neighbour, of weight at most six,
so $\sum_vw_v(z_v^2-\sum_{u\sim v}w_u^2)\le 12\sum_{a<b}w_aw_b\le6D^2$,
while $\sum_vw_v\sum_{u\sim v}w_u^2=\sum_{ab\in F}w_aw_b(w_a+w_b)\le12W$.
Thus $W^2-6qW-12q^3\le0$ and $W\le3q+q\sqrt{9+12q}$. Every edge of $P$ has
its middle pair in the shadow of the core endpoints, and a pair $ab$
accounts for at most $w_aw_b$ edges of $P$, so $q(q-2)/4\le3W$. Dividing
by $q$ gives $q\le38+12\sqrt{12q+9}$, and for $q>38$ squaring gives
$q^2-1804q+148\le0$, which fails for every integer $q\ge1804$.
\end{proof}

For connected $H$ let $\beta=2m-n+1$ be the cycle rank of the incidence
graph over $\mathbb F_2$, and $\Delta=\sum_{\deg v\ge2}(\deg v-2)$.

\begin{lemma}\label{lem:cycle}
If $H$ is connected, $\beta\ge A+2B/3$.
\end{lemma}

\begin{proof}
Group good-leaf hyperedges by their pair $ab$. In a supported group fix a
leafless hyperedge on $ab$ as reference; in an unsupported group (all
hyperedges on $ab$ are good-leaf hyperedges, at least two) fix one of them.
Delete all non-reference good-leaf hyperedges with their leaves. Each
deletion removes two nodes and three edges of the incidence graph without
disconnecting it, so it lowers the cycle rank by one: with $g$ unsupported
groups and reduced hypergraph $R$, $\beta(H)=\beta(R)+A+B-g$.

For an unsupported group $abx,aby$ of size two, the triple $axy$ is
absent, $x$ and $y$ are not joined (the middle hyperedge would contain
$ab$), so $a$ is joined to $x$ or $y$; removing the leaf's hyperedge gives
a two-edge path from $a$ to $b$ avoiding the group. Replacing deleted
hyperedges on this path by their references keeps a genuine two-edge path
(a common reference would be a leafless hyperedge on $ab$). With the
retained reference on $ab$ it forms a cycle of $R$ through three
hyperedges. Project the cycle space of $R$ onto one incidence coordinate
of each of the $g_2$ references of size-two groups. Each chosen cycle
contains its own coordinate and at most three coordinates, so the span has
dimension at least $g_2/3$. Hence $\beta(H)\ge A+\sum(k-1)+g_2/3$ over
unsupported groups of sizes $k$, and $k-1+[k=2]/3\ge2k/3$.
\end{proof}

\begin{lemma}\label{lem:degree}
$\Delta\ge A/2+2B/3$.
\end{lemma}

\begin{proof}
Let $E=abx$ be a good-leaf hyperedge and $F=aby$ another hyperedge on $ab$.
If $a$ and $b$ both have degree two, no pair of the absent triple $axy$ is
joined: a path of length three from $x$ starts with $E$ and continues with
$F$ through $a$ or $b$, after which only $y$ or a vertex of $E$ remains and
no third hyperedge is available; a path between $a$ and $y$ starts with
$E$ or $F$, passes through $b$ and then through the other of $E,F$, and
again stops. So one of $a,b$ has degree at least three. In an unsupported
group $abx,aby$ of size two, if $a$ had degree two, again $axy$ would have
no joined pair: $x$ and $y$ are not joined because a third hyperedge on $ab$
is missing, and a path from $x$ to $a$ would run $x,b$ and end with $aby$,
needing a middle hyperedge from $b$ to the leaf $y$; similarly for $y$.
Larger unsupported groups give degree at least three directly. Charge each supported good leaf $1/2$ to
its endpoints of degree at least three (split equally) and each
unsupported good leaf $1/3$ to each endpoint. A vertex of degree $d\ge4$
receives at most $d/2\le d-2$. A vertex of degree three in at most two
good-leaf hyperedges receives at most one; if all three of its hyperedges
are good-leaf hyperedges, they form one unsupported group of size three
and it receives exactly one. So no vertex receives more than $\deg v-2$.
\end{proof}

\begin{proof}[Proof of the lower bound in Theorem~\ref{thm:main}]
If $H$ is disconnected, Lemma~\ref{lem:basic} gives $m\ge2(n-1)/3$, which
implies~\eqref{eq:lower}. If $H$ is connected, it has no isolated vertex
and counting incidences gives $2n=3m+A+B+q+2t-\Delta$. With
Lemma~\ref{lem:degree} this gives $2n\le3m+A/2+B/3+q+2t$, and
Lemma~\ref{lem:cycle} gives $n\le2m+1-A-2B/3$. Adding twice the first to
the second,
\[
5n\le8m+2q+4t+1 .
\]
Lemma~\ref{lem:pendant} gives $t\le4$, and Lemma~\ref{lem:sqrt} gives
$q\le1803$ and $q\le1+\sqrt{432m+1}$, which yield the two inequalities.
For the third term of the abstract put $x=\sqrt{432m+1}$; the second
inequality becomes $(x+54)^2\ge270n+1891$, so
$x\ge\sqrt{270n+1891}-54>0$, and squaring and substituting
$m=(x^2-1)/432$ gives $m\ge(5n+89)/8-\sqrt{270n+1891}/4$.
\end{proof}

\section{Cones}\label{sec:cones}

\begin{proof}[Proof of Theorem~\ref{thm:cone}]
A Berge-$C_4$ in a cone with root $r$ either avoids $r$ (a $C_4$ in $G$)
or has core $r,a,b,c$ and link edges $ax,ab,bc,cy$ with $x,y\ne b$; then
$x=c$ or $y=a$ gives a paw, $x=y$ a $C_4$ and $x\ne y$ a $P_5$. Conversely
$P_5$, $C_4$ and the paw give Berge-$C_4$s. So $H$ is free if and only if
$G$ contains no $P_5$, $C_4$ or paw, that is, every component is a star,
a double star, a triangle or an isolated vertex. A Berge path whose ends
lie in $\{r\}$ and one component stays in that component.

A component $K_{1,2}$ with leaves $x,y$ fails at the triple $rxy$. A star
centre has no path of length two or three to $r$, so two star components
fail at $ruv$ for their centres. An isolated vertex $z$ requires every pair
of other vertices to be joined; this fails for a star centre and $r$, for
a double-star centre and $r$, and for a second isolated vertex. These
cases give necessity. For sufficiency, each allowed component with $r$ is
saturated: $K_{1,1}$ trivially, $K_{1,t}$ ($t\ge3$) through a third leaf,
$K_3$ because the cone is $K_4^{(3)}$ minus one triple, and $D(p,q)$ by
Lemma~\ref{lem:base} with $Q=\emptyset$. Every vertex of a double star or
triangle has paths of lengths one and two to $r$, every vertex of a star
one of length one, so triples meeting two components are closed through
$r$; with an isolated vertex all pairs in the triangles are joined.

For the sizes, put $N=n-1$. Without an isolated vertex let $\tau$ be the
number of tree components; then $m=N-\tau$, every tree has at least four
vertices except at most one $K_{1,1}$, so $N\ge4\tau-2$, and $\tau=0$
needs $3\mid N$. With an isolated vertex $m=N-1$. Conversely, for $1\le
\tau\le\lfloor(N+2)/4\rfloor$ take $\tau-1$ paths $P_4$ and one double
star on the remaining $N-4(\tau-1)\ge4$ vertices, or, when $N\in\{4\tau-2,
4\tau-1\}$, one $K_{1,1}$, $\tau-2$ paths $P_4$ and a double star on four
or five vertices; $N/3$ triangles give $m=N$. The minimum is
$N-\lfloor(N+2)/4\rfloor=\lceil(3n-5)/4\rceil$.
\end{proof}

Corollary~\ref{cor:interval} follows, since
$\lceil(3n-5)/4\rceil=\lceil(3n-9)/4\rceil+1$ and the smallest value is
realized by Section~\ref{sec:upper}.

\section{The cost of a flexible root}\label{sec:root}

\begin{proof}[Proof of Theorem~\ref{thm:root}]
Let $G$ be the link of $r$ on $N=n-1$ vertices, $a=|E(G)|$, and $b$ the
number of hyperedges avoiding $r$. One-edge paths mean $G$ has no isolated
vertex. Let $t_2$ count components $K_2$ and $t_3$ components $P_3$; every
other component $C$ has at least $3|C|/4$ edges, so $a\ge3N/4-t_2/2-t_3/4$.
A vertex all of whose neighbours in $G$ have degree one, and which lies
only in hyperedges through $r$, has no two-edge path to $r$; so both
vertices of each $K_2$ and the centre of each $P_3$ lie in hyperedges
avoiding $r$, and $3b\ge2t_2+t_3$. Hence $m=a+b\ge3N/4+t_2/6+t_3/12$.
Equality forces $t_2=t_3=0$, $b=0$ and components with exactly $3|C|/4$
edges, that is trees on four vertices; a star $K_{1,3}$ has a centre
without a two-edge path, so every component is $P_4$, which is the block.
\end{proof}

\section{Exact small values}\label{sec:small}

For $n\le2$ the empty hypergraph is vacuously saturated. For $n\ge3$ a
saturated hypergraph is nonempty, and for $n\ge4$ it has at least three
hyperedges, since adding a triple to at most two hyperedges leaves fewer
than four. The witnesses $\{012\}$, $\{012,013,023\}$ and the block show
$s(3)=1$ and $s(4)=s(5)=3$. Cones over the double star $D(1,2)$, over
$K_{1,1}\cup P_4$ and over $K_{1,1}\cup D(1,2)$ give $s(6)\le4$, $s(7)\le4$
and $s(8)\le5$, and Section~\ref{sec:upper} gives the upper bounds for
$9\le n\le13$.

The remaining lower bounds exclude saturated hypergraphs with $m$
hyperedges on $n$ vertices for $m=3$, $n\ge6$; $m=4$, $n\ge8$; $m=5$,
$n\ge10$; $m=6$, $n\ge12$; and $m=7$, $n\ge13$. By
Lemma~\ref{lem:basic}, an isolated vertex forces $3m\ge2(n-1)$, which fails
in each range, and without isolated vertices $n\le3m$. It therefore
suffices to treat hypergraphs without isolated vertices on $n\le3m$
vertices. Each is encoded, after labelling its hyperedges, by the multiset
of incidence masks of its vertices; masks with at least two bits are
listed in nondecreasing order and the remaining incidences are filled by
degree-one vertices. Every hypergraph appears (possibly several times).
The certificate gives, for every profile in the enumeration order, an
exclusion: three vertices in different components, an explicit Berge-$C_4$
of $H$, or an absent triple whose addition leaves the hypergraph free. The
verifier regenerates the enumeration independently, rejects missing or
surplus records, and checks freeness of $H+f$ by testing all four-sets of
hyperedges and all injective choices of cycle vertices. The record counts
are $21$, $261$, $4664$, $107740$ and $10592111$. A sixth certificate
covers all $203895$ profiles with $m=6$, $n=11$; it excludes all but
$1080$, each of which is matched, under an explicit permutation of the
hyperedge labels, to one of the two hypergraphs of Theorem~\ref{thm:small}.
The two are not isomorphic, having degree sequences
$(4,3,2,2,1^7)$ and $(4,3,3,1^8)$.

\section{Formal verification}\label{sec:lean}

The Lean~4 project (Lean~4.35.0, Mathlib
\texttt{56a005855ea17e0b0958f73ec801ae140b0a1679}) defines a hypergraph
on a vertex type $V$ as a finite set of finite subsets of $V$. A
Berge-$C_4$ is given by an injective map $v:\mathrm{Fin}\,4\to V$ and an
injective map $e$ to hyperedges of $H$ with $v_i,v_{i+1}\in e_i$;
\texttt{IsBergeC4Saturated} requires three-uniformity, freeness and a
Berge-$C_4$ after inserting any absent three-element set; and
\texttt{satNumber n} is the least size of a saturated hypergraph on
$\mathrm{Fin}\,n$. The project proves:
\begin{itemize}
\item \texttt{exists\_saturated}: for every $n\ge9$ there is a saturated
hypergraph on $\mathrm{Fin}\,n$ with $(3n-6)/4=\lceil(3n-9)/4\rceil$
hyperedges, and hence \texttt{satNumber\_le\_upper};
\item \texttt{glued\_saturated}: the rooted gluing theorem for an
arbitrary base and block satisfying the hypotheses of
Lemmas~\ref{lem:base} and~\ref{lem:glue}, with the freeness argument
around the root formalized in general;
\item \texttt{satNumber\_small}: $s(n)=0,0,0,1,3,3$ for $n\le5$, and
\texttt{satNumber\_le\_small}: $s(6)\le4$, $s(7)\le4$, $s(8)\le5$.
\end{itemize}
The properties of the explicit bases and the block (uniformity, the
three vertices of degree at least two, paths for every absent triple and
paths to the root) are Boolean searches evaluated by the kernel with
\texttt{decide +kernel}, each with a soundness theorem. No
\texttt{native\_decide} is used, and a module-level audit shows that every
declaration depends only on \texttt{propext}, \texttt{Classical.choice}
and \texttt{Quot.sound}. The lower bound, Theorems~\ref{thm:cone}
and~\ref{thm:root}, Corollary~\ref{cor:interval} and the values
$s(n)$ for $6\le n\le13$ are not formalized.

\section{Related work and scope}

The definitions and the bound $s(n)\le n-2$ are from~\cite{EGKMS}, whose
journal version is~\cite{EGKMS-J}; the inspected statement is that of the
preprint, as the journal text could not be accessed. Linear growth of
$\sat_k(n,\text{Berge-}F)$ for $3\le k\le5$ is due to~\cite{EGMT}, and
gluing at a vertex is used in~\cite{EGKMS}. Wang and Zhang~\cite{WZ} bound the
linear saturation number of Berge-$C_4$, a different problem. The ordinary
path saturation numbers are due to K\'aszonyi and Tuza~\cite{KT}. Searches
in October 2026 of arXiv, journal records and general web indexes did not
find the coefficients $5/8$ or $3/4$, the cone classification or the
values of Theorem~\ref{thm:small}; some recent full texts were
inaccessible, and this does not establish priority.

\section*{Research process and AI disclosure}

This is fully autonomous agentic research. Dev D. Goyal is responsible
for publication. No expert validation or external peer review is
claimed, and the novelty assessment is provisional.

The upper construction was found by autonomous AI research agents in an
earlier session. In a later session the agents proved the lower bound,
the cone classification and spectrum and the root-cost theorem, built and
independently verified the finite certificates and searched the
literature. Their records do not name the system used; the research
contract specifies Astra referees for the internal reviews, which were
AI-assisted reviews of the proofs and not peer review, and the records
state that no Lean formalization was attempted.

Opus 5.5 by Anthropic used with Claude Code assessed the package for this
release. It checked the written proofs and the arithmetic of the lower
bound, compared the statements with~\cite{EGKMS,EGMT,WZ}, rechecked the
construction with its own exhaustive checker, reran the six certificates,
wrote and ran the Lean formalization described in Section~\ref{sec:lean},
and wrote this note. AI checks are not human peer review.

\begin{thebibliography}{9}
\bibitem{EGKMS}
S.~English, N.~Graber, P.~Kirkpatrick, A.~Methuku and E.~C.~Sullivan,
\emph{Saturation of Berge hypergraphs},
\href{https://arxiv.org/abs/1710.03735v1}{arXiv:1710.03735v1} (2017).

\bibitem{EGKMS-J}
S.~English, P.~Gordon, N.~Graber, A.~Methuku and E.~C.~Sullivan,
\emph{Saturation of Berge hypergraphs}, Discrete Mathematics
\textbf{342} (2019), 1738--1761.
\href{https://doi.org/10.1016/j.disc.2019.01.031}{doi:10.1016/j.disc.2019.01.031}.

\bibitem{EGMT}
S.~English, D.~Gerbner, A.~Methuku and M.~Tait,
\emph{Linearity of saturation for Berge hypergraphs}, European Journal of
Combinatorics \textbf{78} (2019), 205--213.
\href{https://arxiv.org/abs/1807.06947}{arXiv:1807.06947}.

\bibitem{KT}
L.~K\'aszonyi and Zs.~Tuza, \emph{Saturated graphs with minimal number of
edges}, Journal of Graph Theory \textbf{10} (1986), 203--210.

\bibitem{WZ}
C.~Wang and J.~Zhang, \emph{Linear saturation numbers of Berge-$C_3$ and
Berge-$C_4$}, \href{https://arxiv.org/abs/2306.06435}{arXiv:2306.06435}
(2023); Applied Mathematics and Computation \textbf{474} (2024).
\end{thebibliography}
\end{document}
